%% RESUMO EM LINGUA ESTRANGEIRA (elemento obrigatorio -- NBR 6028)
%% Versao fiel do resumo.tex, na lingua estrangeira escolhida (geralmente ingles).

\begin{abstract}
Is artificial intelligence already changing the career trajectories of Brazilian
workers? This dissertation investigates whether occupational exposure to
artificial intelligence is associated with changes in labor market transition
patterns following the public release of ChatGPT on November 30, 2022. The study
uses the rotating panel structure of the Brazilian Continuous National Household
Sample Survey (PNAD Contínua) from 2019 to 2025, with an analytical sample of
3,963,359 person-transitions with exposure assigned to the origin occupation,
covering 413 occupational codes and 27 quarters. Exposure is measured using an
occupational artificial intelligence exposure index originally developed for the
U.S. occupational classification and mapped to the Brazilian classification
through a documented crosswalk. The empirical strategy adopts a
difference-in-differences design with continuous treatment, estimated using a
weighted linear probability model with fixed effects for origin occupation and
for state by quarter, and with teleworkability interacted with the period
indicator. After November 2022, the results indicate a general increase in
measured occupational mobility and, beyond this broader movement, a gradient
associated with artificial intelligence exposure: a one-unit increase in the
exposure index is associated with a change between the two periods that is 2.19
percentage points larger in transitions to less exposed occupations and 1.36
percentage points smaller in transitions to more exposed occupations. The most
striking finding, however, concerns employment status transitions: greater
exposure is associated with lower job separation and a lower probability of
moving from formal to informal employment, challenging a simple interpretation
of technological displacement. Taken together, the results are consistent with
an adjustment that shows up in the recomposition of occupational trajectories
rather than in the elimination of jobs, although the design does not isolate the
channel behind this pattern. The estimates should be interpreted as conditional
associations rather than causal effects: the study reports a single estimation
scenario, with no sensitivity battery, and the series used spans a discontinuity
in the measured rate of occupational change that coincides with the change in
the survey's data collection mode in 2020 and 2021. The main contribution of
this study is to provide large-scale Brazilian evidence for a debate still
largely shaped by international research, with implications for public policies
on employment, workforce training, and social protection.
\end{abstract}
